\section{Introduction}\label{sec:intro}

\subsection{Motivation}
\guide{Why $\LTLf$ synthesis matters (planning, robotics, business processes). Why classic
synthesis treats the environment as fully adversarial, and why that is often too pessimistic or
too expensive. One concrete scenario where something is \emph{known} about inputs or outputs.}
Reactive LTL synthesis is a problem in which we are given a specification $\varphi$ in Linear Temporal Logic (LTL) and we want to construct a controller that interacts with every environment such that the resulting trace satisfies the specification.

The TACAS'24 paper of Akshay et al.~\cite{AkshayBCF2024} introduced variable dependencies in LTL synthesis: some output variables are determined by the others, so they can be taken out of the game, which makes the problem smaller.
The definition is made for variables the system sets.
An input dependency would mean that every trace satisfying $\varphi$ fixes the value of an input, even though the environment is free to choose it.
Benchmark specifications rarely constrain their own inputs that way, and when we looked for input dependencies in the benchmarks we found none.

We wanted to generalize the idea of dependencies to input variables as well, and ended up with an even more general idea, in which we solve the synthesis problem with the help of an outside helper that decides the values of some variables.

\subsection{External Knowledge}

Given a specification $\varphi$, we want to be able to incorporate external knowledge about some variables.
We model the knowledge as a transducer that decides the values of these variables from the history of the other variables.
When solving the synthesis problem, we consult the transducer for the values of the variables it controls, which limits the traces to the allowed ones.

The transducer can come from the formula itself, as with output dependencies, or from outside it, as an extra assumption about the environment.

Knowledge buys two things.
First, \textbf{realizability}: some instances are unrealizable when the knowledge is ignored and realizable once it is used.
In the running example (\cref{ex:running}), $\varphi = X[!]\,\top \wedge (o \leftrightarrow X\,k)$ asks the controller to announce the value that $k$ will take at the next step.
If $k$ is a free input, the environment sets it against the announcement; if $k$ is known to echo $r$ one step late, the controller announces $r$ and wins.
Second, \textbf{a smaller product}: the knowledge rules out the traces it does not allow, so the product contains only the states reachable under consistent traces, which can be far fewer than the states of the goal automaton (\texttt{cons-prunes}, \cref{sec:benchmarks}).

Like a fairness assumption, knowledge restricts the traces the controller must handle, but over the same variables it is strictly more expressive than an assumption written in the formula.
When the language of $\Tin$ is expressible in $\LTLf$ by a formula $\psiin$, the problem reduces to standard synthesis of $\psiin \rightarrow \varphi$ (for input knowledge alone; with output knowledge $\psiout$ also appears as a guarantee, \cref{sec:problem}), and knowledge is a different way to say the same thing.
Such instances are the easy case, and the benchmark families whose knowledge has a known $\psiin$ are the ones that can be compared against standard synthesis (\cref{sec:benchmarks}).
But $\LTLf$ expresses only star-free languages, and a transducer need not be star-free.
When the language of $\Tin$ is not star-free, no $\psiin$ exists, so the reduction is unavailable and no $\LTLf$ assumption over the same variables captures the knowledge.
The extension therefore handles inputs that an assumption of the form $\psiin \rightarrow \varphi$ over those variables cannot.
(The qualifier matters: the controller sees the full history, so it could track the state of $\Tin$ itself, and adding that state as fresh known variables would make a local $\psiin$ exist over the larger variable set.)

\subsection{Contributions}
\guide{Bullet list. Problem formulation; the five methods (NFA product, DFA product, three
on-the-fly variants); knowledge extraction from dependencies; implementation on Spot; benchmark
suite and results. Mark which items are done vs.\ in progress.}

\begin{itemize}
  \item Definitions of \begin{itemize}
    \item External knowledge
    \item $\LTLf$ synthesis with external knowledge
    \item $\LTLf$ synthesis with input dependencies
  \end{itemize}
  \item Extraction of knowledge transducers from input and output dependencies
  \item Several benchmark families to test the problem
  \item Implementation of methods to solve the problem, comparisons between them, and a comparison against the state of the art.
\end{itemize}
